\section{Shielded replay SAC and morphology-aware validation}
\label{sec:applied_drl}

The learning and validation protocol uses SAC as the continuous-control backbone because it is a standard off-policy actor--critic method for stochastic policies in continuous action spaces \cite{haarnoja2018soft}. The methodological emphasis lies in the engineering protocol around this backbone: replay warm-starting, action shielding, and morphology-aware checkpoint validation for battery storage dispatch.

\subsection{MDP formulation}

The microgrid scheduling problem is represented as a finite-horizon Markov decision process $(\mathcal{S}, \mathcal{A}, \mathcal{P}, r, \gamma)$ with horizon $T$. The state $s_t \in \mathcal{S} \subset \mathbb{R}^{36}$ aggregates time features (hour-of-day, day-of-week encoded cyclically), current battery SOC, normalized active and reactive load, photovoltaic output, electricity price, and compact network-stress indicators including voltage-band status and feeder-loading severity. The action $a_t \in \mathcal{A} = [-1, 1]$ is the normalized BESS active-power command, where $+1$ corresponds to maximum discharge and $-1$ to maximum charge. After denormalization and physical scaling, the simulator applies battery-side clipping and optional shield correction before writing the achieved terminal power $P_{\mathrm{bess}}^{\mathrm{ach}}(t)$ into the network model. The transition function $\mathcal{P}(s_{t+1} \mid s_t, a_t)$ is deterministic conditional on the prescribed exogenous load, PV, and price profiles and the battery SOC update law in Eq.~\eqref{eq:soc_update_morphology}.

Figure~\ref{fig:sac_architecture} illustrates the shielded SAC framework for battery dispatch control, showing the integration of replay warm-starting, action shielding, and network-settlement feedback.

\begin{figure}[htbp]
\centering
\includegraphics[width=0.95\columnwidth]{sac_algorithm.pdf}
\caption{Shielded SAC architecture for battery dispatch. The framework integrates replay warm-starting (oracle-seeded buffer), action shielding (battery feasibility enforcement), and network-settlement feedback (pandapower power flow) with morphology-aware checkpoint validation.}
\label{fig:sac_architecture}
\end{figure}

The training reward is cost-led but not identical to the final evaluation objective. It is constructed as
\begin{equation}
\begin{aligned}
r_t ={}& -c_{\mathrm{grid}}(t)
- \lambda_{\mathrm{soc}} \psi_{\mathrm{soc}}(s_t)
- \lambda_{\mathrm{net}} \psi_{\mathrm{net}}(s_t) \\
&- \lambda_{\mathrm{inf}} \psi_{\mathrm{inf}}(a_t)
- \lambda_{\mathrm{shape}} \psi_{\mathrm{shape}}(s_t),
\end{aligned}
\label{eq:training_reward}
\end{equation}
where $c_{\mathrm{grid}}(t)$ is the grid energy cost, $\psi_{\mathrm{soc}}$ penalizes SOC boundary saturation and deviation from mid-band residence, $\psi_{\mathrm{net}}$ penalizes voltage violations and line overloading, $\psi_{\mathrm{inf}}$ penalizes infeasible raw battery actions before shield correction, and $\psi_{\mathrm{shape}}$ includes auxiliary shaping terms for terminal inventory consistency and discharge reserve during high-price periods. The evaluation objective reported in the results uses logged operating cost and terminal SOC penalty rather than the full clipped training reward in Eq.~\eqref{eq:training_reward}. This distinction matters because a controller can improve training reward through shaping terms or shield corrections without producing an operator-interpretable evaluation trajectory. The main evaluation reports six P4-selected seeds from the trained candidate ensemble as a focused basis for evaluating morphology and cross-fidelity behavior. Each seed requires approximately 18--24 wall-clock hours on a workstation with an Intel Xeon Gold 6248R CPU and 128\,GB RAM, so ensemble training is computationally material. Modern DRL practice recommends 20--30 seeds for formal inference \cite{agarwal2021statistical}; the present scope is therefore treated as an engineering evaluation rather than a formal statistical comparison.

\subsection{Replay warm-starting and shielding}

The replay buffer is initialized with $N_{\mathrm{seed}} = 10{,}000$ transitions (approximately 10\% of the final buffer capacity of 100{,}000) generated by a linear-programming oracle with perfect foresight over 30-day training windows. The oracle schedule is replayed through the same network-settlement interface and shield layer used during online training, so the seeded trajectories respect battery feasibility and terminal-inventory targets while exposing the actor and critic to economically informative storage behavior before online SAC updates dominate the buffer. Replay warm-starting improves early learning stability, but it can also create over-dependence on oracle or shield corrections if the online policy fails to explore beyond the seeded region. For this reason, oracle-teacher activation, shield intervention, and correction magnitude are reported as evaluation diagnostics rather than hidden implementation details.

The action shield receives the actor command and the current battery state, then returns an achieved action that respects battery feasibility and terminal-inventory targets. Let $a_{\theta}(t)$ be the actor action and $a_{\mathrm{sh}}(t)$ the shielded action. The shield correction magnitude is
\begin{equation}
\Delta a_{\mathrm{sh}}(t)=\left|a_{\mathrm{sh}}(t)-a_{\theta}(t)\right|.
\label{eq:shield_delta}
\end{equation}
The inventory-teacher action gap, $g_{\mathrm{IT}}$, is the trajectory-averaged absolute distance between the actor's normalized command and the normalized command produced by the inventory-teacher rule:
\begin{equation}
g_{\mathrm{IT}}=\frac{1}{T}\sum_{t=1}^{T}\left|a_{\theta}(t)-a_{\mathrm{IT}}(t)\right|,
\label{eq:inventory_teacher_gap}
\end{equation}
where $a_{\mathrm{IT}}(t)$ is the normalized reference command produced by the inventory-teacher rule at the same state. Because both commands use the normalized action scale $[-1,1]$, $g_{\mathrm{IT}}$ is a dimensionless action-space discrepancy rather than an additional cost term. The mean shield correction, activation fractions, and inventory-teacher gap are aggregated over each evaluation window. A low correction magnitude suggests greater policy internalization, whereas high correction magnitudes indicate that the executed controller remains a shielded protocol rather than an autonomous learned policy. A nonzero $g_{\mathrm{IT}}$ is neither a shield activation nor a corrective action; it measures separation between the learned action and the teacher reference before trajectory-level averaging.

Table~\ref{tab:action_diagnostic_definitions} separates the action and intervention quantities used in the audit. The distinction is important because infeasible raw actions, shield corrections, and teacher-action gaps are related but not equivalent diagnostics.

\begin{table}[!t]
\centering
\footnotesize
\setlength{\tabcolsep}{3.5pt}
\renewcommand{\arraystretch}{1.08}
\caption{Action and intervention diagnostics used by the validation protocol.}
\label{tab:action_diagnostic_definitions}
\begin{tabularx}{\columnwidth}{lX}
\toprule
Quantity & Interpretation \\
\midrule
$a_{\theta}(t)$ & Raw normalized actor command before battery feasibility checks. \\
$a_{\mathrm{sh}}(t)$ & Executed normalized action after clipping and optional shield correction. \\
$a_{\mathrm{IT}}(t)$ & Inventory-teacher reference action at the same state. \\
Infeasible-action dwell & Fraction of raw actor commands outside the battery feasibility envelope. \\
Shield activation & Fraction of steps with non-negligible executed-action correction. \\
$\Delta a_{\mathrm{sh}}(t)$ & Magnitude of the actor-to-executed-action correction. \\
$g_{\mathrm{IT}}$ & Mean normalized separation between actor and teacher actions. \\
\bottomrule
\end{tabularx}
\end{table}

\emph{Remark on safe RL formulation.} The method uses action shielding rather than policy-constrained formulations such as constrained policy optimization \cite{achiam2017cpo} or control-barrier functions. This choice reflects the engineering focus of the study. The shield is a transparent safety layer whose activation can be measured and reported, whereas policy-constrained methods internalize constraints during training but provide limited post-hoc visibility into whether the learned policy respects constraints autonomously or relies on implicit Lagrangian correction. By reporting shield activation explicitly, the protocol avoids conflating shielded execution with policy internalization, a distinction used by the morphology-aware validation framework.

\subsection{Morphology-aware checkpoint validation}

The validation score used by the BalancedGate protocol augments objective cost with inventory and behavior terms. In compact form, the validation objective can be written as
\begin{equation}
\begin{aligned}
J_{\mathrm{val}} = J_{\mathrm{obj}} &+ \lambda_{\mathrm{term}}D_{\mathrm{term}} + \lambda_{\mathrm{bd}}D_{\mathrm{bd}} + \lambda_{\mathrm{inf}}D_{\mathrm{inf}} \\
&+ \lambda_{\mathrm{peak}}D_{\mathrm{peak}} + \lambda_{\mathrm{sh}}D_{\mathrm{sh}},
\end{aligned}
\label{eq:morphology_validation}
\end{equation}
where $J_{\mathrm{obj}}$ is the validation objective cost, $D_{\mathrm{term}}$ is terminal SOC deviation or penalty, $D_{\mathrm{bd}}$ is boundary parking, $D_{\mathrm{inf}}$ is infeasible-action dwell, $D_{\mathrm{peak}}$ measures insufficient high-price discharge behavior, and $D_{\mathrm{sh}}$ summarizes shield or teacher-action dependence. The exact implementation uses logged diagnostic fields, but Eq.~\eqref{eq:morphology_validation} captures the design principle: checkpoint selection is not driven by objective value alone. Each $\lambda_k D_k$ term may aggregate multiple logged fields; for example, $D_{\mathrm{sh}}$ combines mean shield correction, material shield activation, strong shield activation, and teacher-gap penalties when available, while $D_{\mathrm{peak}}$ incorporates both peak-reserve headroom shortfall and peak-discharge behavior penalties. The weights and hard-gate thresholds are logged in the released run summaries and configuration fields. These weights are case-dependent engineering penalties chosen for the IEEE33 network-stress case. They are not universal constants, and the contribution of this paper is the validation methodology itself rather than any specific weight configuration.

The penalty weights were selected by a heuristic calibration procedure rather than by test-year optimization. First, diagnostic terms were expressed as fractions, normalized action gaps, or objective-equivalent terminal penalties so that their numerical ranges were comparable on validation windows. Second, infeasible-action and shield-related penalties were set to dominate small objective-cost differences once operator-facing gates were violated, while allowing cost to rank checkpoints inside the acceptable morphology region. Third, the shield term was bracketed through the P1--P4 protocol variants: P1 applies a stricter shield burden, P2 removes the strict shield gate, and P4 uses the intermediate setting reported in the main evaluation. For transfer to another feeder or market setting, the same procedure would start from the no-storage and reference-objective scale, choose operator-acceptable morphology thresholds, and then tune $\lambda_k$ on validation windows before any held-out-year evaluation.

The P4 gates used in the case study are screening thresholds rather than universal operating standards. The mid-band and infeasible-action cutoffs are deliberately permissive so that the protocol can retain candidates for detailed cross-year and cross-fidelity audit. Stricter operator standards can be imposed by tightening the same gates before deployment-oriented selection.

\subsection{Protocol variants}

Four protocol variants characterize the value--morphology tradeoff. All variants share the same SAC backbone, replay warm-starting, and shield; they differ only in the validation metric used for checkpoint selection. P1 applies a strict validation gate that assigns large penalties to boundary dwell and shield activation, producing conservative checkpoints with low risk but also low value recovery. P2 removes the shield-activation gate while retaining the morphology-aware weighted sum, permitting higher value recovery at the cost of greater intervention dependence. P3 uses a coarser checkpoint interval (every 25\% of training rather than every 5\%) with the same weighted validation metric; it exposes the risk of selecting intervention-dependent checkpoints when validation checks are less frequent. P4 BalancedGate re-enables a softer gate on shield metrics and is the proposed reporting variant.

The final selection logic combines value, morphology, and protocol-dependence status. Positive value recovery with clean or improved-but-fragile morphology is retained for engineering interpretation. Positive value with problematic morphology is reported as behavior-fragile, and negative value with poor morphology is not retained as a dispatch candidate. This logic is stricter than reporting scalar reward alone and reflects engineering expectations for storage control.
